\scanpage{23}
\uline{Contour \inserted{(comb.)} connexe} \inserted{$C$} définit

\textcircled{1} L'ens $\omega=\omega(C)$ à deux \unsure{éléments|élémts} de \unsure{ses|ces} deux orientations.

À partir de là, on peut déjà écrire (à isom can. près) le groupe des translations du contour, \unsure{par}
\[
 G^+=\operatorname{Aut}^+(C)\xleftarrow{\sim}\omega\wedge_{\{\pm1\}}\mathbb Z_n
\]
(où $\{\pm1\}$ opère sur $\mathbb Z_n$ par \unsure{homothéties|homotheties})

\textcircled{2} Un torseur sous $G^+$

\deleted{\uline{\unsure{Exercice|Exercices}}} Et \unsure{inversmt|inversemt|inversement}, le couple formé par $(\omega,T)$, où $\omega$ un \deleted{\unsure{tors|ens}} à $2$ él et $T$ un torseur sous $G^+\mathrel{\unsure{=|\simeq|\leftrightarrows}}\omega\wedge_{\{\pm1\unsure{,}\}}\mathbb Z_n$, définit un
% archive sheet number: 22
